\vspace{0.5em}\noindent\textbf{BUILDING A SCALABLE REAL-TIME CHAT SYSTEM WITH SOCKET.IO,
REDIS, AND
DYNAMODB}\par

Realtime chat appears in many systems such as e-commerce, recruitment
platforms, customer support, social networks, and collaboration tools.

At a basic level, a single server can receive and send messages.
However, when the number of users grows, the system needs to answer more
difficult questions:

\begin{itemize}
\tightlist
\item
  How can messages appear immediately?
\item
  What happens when the sender and receiver are connected to different
  servers?
\item
  Where should messages be stored so they are not lost when a server
  restarts?
\item
  How can the number of servers increase without interrupting users?
\end{itemize}

A common architecture for solving these problems combines Socket.IO,
Redis, and Amazon DynamoDB.

\vspace{0.5em}\noindent\textbf{1. Why REST API is not enough for realtime
chat}\par

REST API follows a request-response model:

\begin{Shaded}
\begin{Highlighting}[]
\NormalTok{Client sends request}
\NormalTok{v}
\NormalTok{Server processes}
\NormalTok{v}
\NormalTok{Server returns response}
\end{Highlighting}
\end{Shaded}

If REST API is used for chat, the frontend usually needs to keep calling
an API to check for new messages:

\begin{Shaded}
\begin{Highlighting}[]
\NormalTok{GET /messages}
\NormalTok{Wait a few seconds}
\NormalTok{GET /messages}
\NormalTok{Wait a few seconds}
\NormalTok{GET /messages}
\end{Highlighting}
\end{Shaded}

This approach is called polling.

Polling is simple, but it has several limitations:

\begin{itemize}
\tightlist
\item
  Message latency depends on the polling interval.
\item
  The client sends requests even when there are no new messages.
\item
  The backend has to process many unnecessary requests.
\item
  System load increases quickly as the number of users grows.
\end{itemize}

For systems that require near-instant responses, WebSocket or a realtime
library such as Socket.IO is usually more suitable.

\vspace{0.5em}\noindent\textbf{2. Socket.IO for realtime
communication}\par

Socket.IO allows a two-way connection to be maintained between the
client and the server.

After the connection is established, the client and server can send
events to each other without creating a new HTTP request for every
exchange.

A simple message flow can be understood as:

\begin{Shaded}
\begin{Highlighting}[]
\NormalTok{User A sends a message}
\NormalTok{v}
\NormalTok{Socket.IO server receives the event}
\NormalTok{v}
\NormalTok{Server processes the message}
\NormalTok{v}
\NormalTok{Server sends a new event}
\NormalTok{v}
\NormalTok{User B receives the message}
\end{Highlighting}
\end{Shaded}

Common events include:

\begin{itemize}
\tightlist
\item
  \texttt{conversation:join}
\item
  \texttt{message:send}
\item
  \texttt{message:new}
\item
  \texttt{user:typing}
\item
  \texttt{message:read}
\end{itemize}

Socket.IO also supports:

\begin{itemize}
\tightlist
\item
  Automatic reconnection when the network is interrupted.
\item
  Grouping users into rooms.
\item
  Sending events to one user or a group of users.
\item
  Fallback when WebSocket is not available.
\end{itemize}

However, Socket.IO works best when the system has only one server. When
the application runs multiple servers, an additional synchronization
mechanism is required.

\vspace{0.5em}\noindent\textbf{3. The issue when the chat service runs multiple
instances}\par

With one chat server, both users connect to the same server:

\begin{Shaded}
\begin{Highlighting}[]
\NormalTok{User A {-}+}
\NormalTok{        +{-}{-} Chat server}
\NormalTok{User B {-}+}
\end{Highlighting}
\end{Shaded}

When User A sends a message, the server can send it directly to User B's
connection.

As traffic grows, the system may need multiple servers:

\begin{Shaded}
\begin{Highlighting}[]
\NormalTok{Chat server 1}
\NormalTok{Chat server 2}
\NormalTok{Chat server 3}
\end{Highlighting}
\end{Shaded}

A load balancer may distribute users like this:

\begin{Shaded}
\begin{Highlighting}[]
\NormalTok{User A $\textbackslash{}rightarrow$ Chat server 1}
\NormalTok{User B $\textbackslash{}rightarrow$ Chat server 2}
\end{Highlighting}
\end{Shaded}

At this point, Chat server 1 does not directly manage User B's
connection.

Without communication between servers, User A's message may be saved
successfully, but User B may not receive a realtime notification.

This is the problem Redis can solve.

\vspace{0.5em}\noindent\textbf{4. Redis synchronizes events between chat
servers}\par

Redis is a high-speed in-memory data store. In realtime chat
architecture, Redis is commonly used for publish/subscribe.

The flow works like this:

\begin{Shaded}
\begin{Highlighting}[]
\NormalTok{User A}
\NormalTok{v}
\NormalTok{Chat server 1}
\NormalTok{v Publish}
\NormalTok{Redis}
\NormalTok{v Subscribe}
\NormalTok{Chat server 2}
\NormalTok{v}
\NormalTok{User B}
\end{Highlighting}
\end{Shaded}

When Chat server 1 receives a message:

\begin{itemize}
\tightlist
\item
  The server processes and stores the message.
\item
  The server publishes an event to Redis.
\item
  Other chat servers receive the event.
\item
  The server that holds the receiver's connection sends the message to
  the client.
\end{itemize}

Redis works like a shared broadcast channel between chat servers.

Redis is suitable for this role because:

\begin{itemize}
\tightlist
\item
  It has high processing speed.
\item
  It supports publish/subscribe.
\item
  It is suitable for short-lived data and events.
\item
  It allows servers to operate independently while still exchanging
  events.
\end{itemize}

Redis does not have to be the long-term storage location for chat
history. That task should be handled by a more durable database.

\vspace{0.5em}\noindent\textbf{5. DynamoDB stores chat
data}\par

Amazon DynamoDB is a managed NoSQL database service from AWS.

Chat data often has these characteristics:

\begin{itemize}
\tightlist
\item
  The number of messages continuously increases.
\item
  Write throughput can be high.
\item
  Messages are usually queried by conversation.
\item
  Complex joins are rarely needed.
\item
  Usage can spike suddenly.
\end{itemize}

DynamoDB fits this model because it scales well and does not require
users to manage database servers.

A simple message table design can use:

\begin{itemize}
\tightlist
\item
  Partition key: \texttt{conversation\_\allowbreak{}id}
\item
  Sort key: \texttt{timestamp} or \texttt{message\_\allowbreak{}id}
\end{itemize}

In this design:

\begin{itemize}
\tightlist
\item
  \texttt{conversation\_\allowbreak{}id} groups messages from the same conversation.
\item
  \texttt{timestamp} sorts messages by time.
\end{itemize}

When chat history is needed, the system queries records with the same
\texttt{conversation\_\allowbreak{}id}.

The roles of Redis and DynamoDB can be separated clearly:

\begin{Shaded}
\begin{Highlighting}[]
\NormalTok{Redis}
\NormalTok{$\textbackslash{}rightarrow$ Transmits realtime events between servers}

\NormalTok{DynamoDB}
\NormalTok{$\textbackslash{}rightarrow$ Stores long{-}term chat history}
\end{Highlighting}
\end{Shaded}

\vspace{0.5em}\noindent\textbf{6. Complete message sending
flow}\par

A message can pass through the system in these steps:

\begin{enumerate}
\def\labelenumi{\arabic{enumi}.}
\tightlist
\item
  The user enters a message.
\item
  The client sends an event through Socket.IO.
\item
  The chat server authenticates the user.
\item
  The server checks whether the user can access the conversation.
\item
  The message is stored in DynamoDB.
\item
  The server publishes an event through Redis.
\item
  The chat server holding the receiver's connection receives the event.
\item
  The server sends the new message to the receiver's client.
\end{enumerate}

General architecture:

\begin{Shaded}
\begin{Highlighting}[]
\NormalTok{Client A}
\NormalTok{|}
\NormalTok{| Socket.IO}
\NormalTok{v}
\NormalTok{Chat server 1}
\NormalTok{|}
\NormalTok{+{-}{-}{-}{-} Store data {-}{-}{-}{-}{-}\textgreater{} DynamoDB}
\NormalTok{|}
\NormalTok{+{-}{-}{-}{-} Publish event {-}{-}\textgreater{} Redis}
\NormalTok{                         |}
\NormalTok{                         v}
\NormalTok{                  Chat server 2}
\NormalTok{                         |}
\NormalTok{                         v}
\NormalTok{                      Client B}
\end{Highlighting}
\end{Shaded}

This architecture achieves two goals:

\begin{itemize}
\tightlist
\item
  Realtime responses through Socket.IO and Redis.
\item
  Durable message history through DynamoDB.
\end{itemize}

\vspace{0.5em}\noindent\textbf{7. How Kubernetes supports scaling chat
services}\par

When the application is packaged as containers, Kubernetes can run
multiple replicas of the chat service.

\begin{Shaded}
\begin{Highlighting}[]
\NormalTok{Chat service}
\NormalTok{+{-}{-} Pod 1}
\NormalTok{+{-}{-} Pod 2}
\NormalTok{+{-}{-} Pod 3}
\NormalTok{+{-}{-} Pod N}
\end{Highlighting}
\end{Shaded}

A Kubernetes Service or load balancer distributes connections to the
pods.

Kubernetes also supports:

\begin{itemize}
\tightlist
\item
  Restarting containers automatically when failures occur.
\item
  Rolling updates for new versions.
\item
  Scaling the number of pods up or down.
\item
  Health checks through liveness and readiness probes.
\item
  Traffic distribution across multiple pods.
\item
  Horizontal Pod Autoscaler based on CPU, memory, or custom metrics.
\end{itemize}

However, Kubernetes only runs and manages multiple instances. Realtime
synchronization between instances still requires Redis or a suitable
message broker.

\vspace{0.5em}\noindent\textbf{8. Sticky session and Redis solve different
problems}\par

When Socket.IO uses HTTP long polling, a client can create multiple
consecutive requests in the same session.

If these requests are routed to different servers, the connection
session may fail.

Sticky session keeps requests from the same client on the same server:

\begin{Shaded}
\begin{Highlighting}[]
\NormalTok{Client A $\textbackslash{}rightarrow$ Chat server 1}
\NormalTok{Client A $\textbackslash{}rightarrow$ Chat server 1}
\NormalTok{Client A $\textbackslash{}rightarrow$ Chat server 1}
\end{Highlighting}
\end{Shaded}

However, sticky session does not replace Redis.

The difference is:

\begin{itemize}
\tightlist
\item
  Sticky session keeps one client's connection stable on one server.
\item
  Redis transmits events between different servers.
\end{itemize}

Even when sticky session is enabled, the sender and receiver may still
be connected to different servers.

\vspace{0.5em}\noindent\textbf{9. Deployment on AWS}\par

On AWS, the architecture can use:

\begin{Shaded}
\begin{Highlighting}[]
\NormalTok{Application Load Balancer}
\NormalTok{v}
\NormalTok{Amazon EKS or container service}
\NormalTok{v}
\NormalTok{Multiple chat service instances}
\NormalTok{+{-}{-} Amazon DynamoDB}
\NormalTok{+{-}{-} Amazon ElastiCache/Valkey}
\end{Highlighting}
\end{Shaded}

Service roles:

\begin{itemize}
\tightlist
\item
  Amazon EKS: runs and manages chat service containers.
\item
  Application Load Balancer: receives and distributes user connections.
\item
  Amazon DynamoDB: stores chat data.
\item
  Amazon ElastiCache or Valkey: provides a Redis-compatible service.
\item
  Amazon CloudWatch: stores logs, metrics, and alarms.
\end{itemize}

For small projects or test environments, Redis and DynamoDB Local can be
run with Docker before using real AWS services.

\vspace{0.5em}\noindent\textbf{10. Common mistakes when building realtime
chat}\par

\vspace{0.5em}\noindent\textbf{Storing messages only in server
memory}\par

If the server restarts, all data may be lost.

\begin{Shaded}
\begin{Highlighting}[]
\NormalTok{Server restart $\textbackslash{}rightarrow$ Message history is lost}
\end{Highlighting}
\end{Shaded}

Messages should be stored in a durable database.

\vspace{0.5em}\noindent\textbf{Running multiple servers without a Redis
adapter}\par

Each server only knows the clients directly connected to it. Users on
another server may not receive events.

\vspace{0.5em}\noindent\textbf{Using Redis as the only storage
layer}\par

Redis is suitable for cache and event transmission, but it is not always
suitable as the only source of chat history.

\vspace{0.5em}\noindent\textbf{Not authenticating Socket.IO
connections}\par

WebSocket or Socket.IO connections still need to be authenticated with
JWT, session, or another appropriate mechanism.

\vspace{0.5em}\noindent\textbf{Not checking conversation
access}\par

A logged-in user is not automatically allowed to join every chat room.

\vspace{0.5em}\noindent\textbf{Not handling duplicate
messages}\par

When the network is unstable, the client may resend the same message. A
client-generated message ID or idempotency key can help reduce duplicate
records.

\vspace{0.5em}\noindent\textbf{Not limiting message size and sending
frequency}\par

The system should include:

\begin{itemize}
\tightlist
\item
  Message length limits.
\item
  Rate limiting.
\item
  Spam prevention.
\item
  File type validation if attachments are supported.
\end{itemize}

\vspace{0.5em}\noindent\textbf{11. Cost optimization ideas}\par

Realtime chat on AWS can generate cost from compute, load balancers,
databases, Redis, and logs.

Some optimization ideas:

\begin{itemize}
\tightlist
\item
  Use DynamoDB On-Demand when traffic is low or unpredictable.
\item
  Choose a suitable ElastiCache node instead of over-provisioning.
\item
  Use Auto Scaling to avoid running too many instances during low
  traffic.
\item
  Set retention periods for CloudWatch Logs instead of keeping logs
  forever.
\item
  Use DynamoDB TTL for temporary data when the business rules allow it.
\item
  Delete Load Balancers, ElastiCache clusters, or test environments
  after labs are completed.
\item
  Run the system with Docker and DynamoDB Local during development.
\end{itemize}

Managed services reduce operational work, but resources still need to be
monitored to avoid unexpected cost.

\vspace{0.5em}\noindent\textbf{12. Reference checklist}\par

\begin{itemize}
\tightlist
\item
  Use Socket.IO or WebSocket for realtime communication.
\item
  Use a Redis adapter when running multiple chat servers.
\item
  Store message history in a durable database.
\item
  Design an appropriate DynamoDB partition key.
\item
  Do not store important state only in server memory.
\item
  Configure sticky session if using long polling.
\item
  Authenticate realtime connections.
\item
  Check conversation access permission.
\item
  Add rate limiting and message size limits.
\item
  Configure liveness and readiness probes.
\item
  Monitor connection count, message rate, latency, and error rate.
\item
  Configure Auto Scaling and resource limits properly.
\item
  Check AWS resource cost after testing.
\end{itemize}

The most important lesson is that realtime chat is not only about
setting up a WebSocket connection.

When the system starts to scale, three responsibilities should be
separated clearly:

\begin{Shaded}
\begin{Highlighting}[]
\NormalTok{Socket.IO}
\NormalTok{$\textbackslash{}rightarrow$ Realtime communication between client and server}

\NormalTok{Redis}
\NormalTok{$\textbackslash{}rightarrow$ Event synchronization between multiple servers}

\NormalTok{DynamoDB}
\NormalTok{$\textbackslash{}rightarrow$ Long{-}term chat data storage}
\end{Highlighting}
\end{Shaded}

Choosing the right tool for each responsibility helps the system respond
quickly while remaining scalable as the number of users increases.

\texttt{\#AWS} \texttt{\#SocketIO} \texttt{\#Redis} \texttt{\#DynamoDB}
\texttt{\#AmazonEKS} \texttt{\#RealtimeChat} \texttt{\#NodeJS}
\texttt{\#Kubernetes} \texttt{\#CloudComputing} \texttt{\#AWSStudyGroup}
